\subsection{乘性子集的更换}

设 A 是环，S 是 A 的乘性子集，M 是 A-模。若 T 是 A 的包含 S 的乘性子集，则由第 2 节命题 5 可知，存在唯一的 \(S^{-1}A\)-线性映射 \(i_{M}^{T,S}:S^{-1}M\to T^{-1}M\)，使得 \(i_{M}^{T}=i_{M}^{T,S}\circ i_{M}^{S}\)；映射 \(i_{M}^{T,S}\) 把 \(S^{-1}M\) 中的元素 m/s 映为 \(T^{-1}M\) 中的元素 m/s。容易验证 \(i_{M}^{T,S}=i_{A}^{T,S}\otimes l_{M}\)。若 U 是 A 的第三个乘性子集且 \(T\subset U\)，则 \(i_{M}^{U,S}=i_{M}^{U,T}\circ i_{M}^{T,S}\)；此外，若 u: \(M\to N\) 是 A-模同态，则图表

\[
\begin{array}{c} \mathrm{S} ^ {- 1} \mathrm{M} \xrightarrow {\mathrm{S} - 1 _ {u}} \mathrm{S} ^ {- 1} \mathrm{N} \\ i _ {\mathrm{M}} ^ {\mathrm{T}, \mathrm{S}} \Biggl \downarrow \qquad \qquad \qquad \qquad \qquad \qquad \qquad \Biggl \downarrow i _ {\mathrm{N}} ^ {\mathrm{T}, \mathrm{S}} \\ \mathrm{T} ^ {- 1} \mathrm{M} \xrightarrow [ \mathrm{T} - 1 _ {u} ]{} \mathrm{T} ^ {- 1} \mathrm{N} \end{array}
\]

是交换的。

命题 7. 设 A 是环，S、T 是 A 的两个乘性子集。令 \(\mathbf{T}' = i(\mathbf{T}_{\mathbf{A}}^{\mathbf{s}})\)。

\begin{enumerate}
\def\labelenumi{(\roman{enumi})}
\tightlist
\item
  存在唯一的同构 \(j\)（从环 (ST) \(^{-1}\) A 到环 \(\mathbf{T}^{\prime -1}(\mathbf{S}^{-1}\mathbf{A})\)），使得图表

\[
\begin{array}{c} \mathrm{A} \xrightarrow {i _ {\mathrm{A}} ^ {\mathrm{S}}} \mathrm{S} ^ {- 1} \mathrm{A} \\ i _ {\mathrm{M}} ^ {\mathrm{ST}} \Big \downarrow \qquad \qquad \qquad \qquad \qquad \qquad \qquad \Big \downarrow i _ {\mathrm{S} - 1 _ {\mathrm{A}}} ^ {\mathrm{T} ^ {\prime}} \\ (\mathrm{ST}) ^ {- 1} \mathrm{A} \xrightarrow [ j ]{} \mathrm{T} ^ {\prime - 1} (\mathrm{S} ^ {- 1} \mathrm{A}) \end{array}
\]

是交换的。

\item
  设 M 是 A-模。存在一个 \((\mathrm{ST})^{-1}\mathrm{A}\)-同构 k，从 \((\mathrm{ST})^{-1}\mathrm{A}\)-模 \((\mathrm{ST})^{-1}\mathrm{M}\) 到 \(\mathrm{T}^{\prime-1}(\mathrm{S}^{-1}\mathrm{A})\)-模 \(\mathrm{T}^{\prime-1}(\mathrm{S}^{-1}\mathrm{M})\)，使得图表
\end{enumerate}

\[
\begin{array}{c} \mathrm{M} \qquad \xrightarrow {i _ {\mathrm{M}} ^ {\mathrm{S}}} \mathrm{S} ^ {- 1} \mathrm{M} \\ i _ {\mathrm{A}} ^ {\mathrm{ST}} \Big \downarrow \qquad \qquad \qquad \qquad \qquad \Big \downarrow i _ {\mathrm{S}} ^ {\mathrm{T} ^ {\prime} - 1 _ {\mathrm{M}}} \\ (\mathrm{ST}) ^ {- 1} \mathrm{M} \xrightarrow [ k ]{} \mathrm{T} ^ {\prime - 1} (\mathrm{S} ^ {- 1} \mathrm{M}) \end{array}
\]

是交换的。

\begin{enumerate}
\def\labelenumi{(\roman{enumi})}
\tightlist
\item
  我们利用 \((\mathrm{ST})^{-1}\mathrm{A}\) 作为一个泛映射问题的解这一定义。设 B 是环，f 是 A 到 B 的同态，使得 \(f(\mathrm{ST})\) 由可逆元组成。由于 \(f(\mathrm{S})\) 从而也由可逆元组成，故存在唯一的同态 \(f': S^{-1}A \to B\) 使得 \(f = f' \circ i_{A}^{S}\)（第 1 节，命题 1）。对所有 \(t \in T\)，\(f'(i_{\mathbf{A}}^{\mathrm{S}}(t)) = f(t)\) 依假设在 B 中可逆，因而 \(f'(\mathrm{T}')\) 由可逆元组成；于是由第 1 节命题 1，存在唯一的同态 \(f''\)（从 \(T'^{-1}(S^{-1}A)\) 到 B）使得 \(f' = f'' \circ i_{S^{-1}A}^{T'}\)，从而 \(f = f'' \circ u\)，其中令 \(u = i_{S^{-1}A}^{T'} \circ i_{A}^{S}\)。

此外，若 \(f_1'' \colon \mathrm{T}^{-1}(\mathrm{S}^{-1}\mathrm{A}) \to \mathrm{B}\) 是使 \(f_1'' \circ u = f\) 成立的第二个同态，则 \((f_1'' \circ i_{\mathrm{S}^{-1}\mathrm{A}}^{\mathrm{T}'} \circ i_{\mathrm{A}}^{\mathrm{S}} = (f'' \circ i_{\mathrm{S}^{-1}\mathrm{A}}^{\mathrm{T}'} \circ i_{\mathrm{A}}^{\mathrm{S}} \text{，由此 } f_1'' \circ i_{\mathrm{S}^{-1}\mathrm{A}}^{\mathrm{T}'} = f'' \circ i_{\mathrm{S}^{-1}\mathrm{A}}^{\mathrm{T}'} \text{ 从而 } f_1'' = f''\)。

由于 ST 的元素经 \(u\) 在 \(\mathbf{T}^{-1}(\mathbf{S}^{-1}\mathbf{A})\) 中的像都可逆，有序对 \((\mathbf{T}^{-1}(\mathbf{S}^{-1}\mathbf{A}), u)\) 是第 1 节中所考虑的（关于 A 与 ST 的）泛映射问题的解。这就表明了 \(j\) 的存在性与唯一性。

\item
  证明与 (i) 完全类似，只是此时要用第 2 节命题 3，留给读者。
\end{enumerate}

命题 8. 设 A 是环，S、T 是 A 的两个乘性子集且 \(\mathbf{S} \subset \mathbf{T}\)。下列性质等价：

\begin{enumerate}
\def\labelenumi{(\alph{enumi})}
\item
  同态 \(i_{\mathbf{A}}^{\mathrm{T},\mathrm{S}}: \mathrm{S}^{-1}\mathrm{A} \to \mathrm{T}^{-1}\mathrm{A}\) 是双射。
\item
  对每个 A-模 M，同态 \(i_{\mathbf{A}}^{\mathrm{T},\mathrm{S}}: \mathrm{S}^{-1}\mathrm{M} \to \mathrm{T}^{-1}\mathrm{M}\) 是双射。
\item
  对所有 \(t \in \mathbf{T}\)，存在 \(a \in \mathbf{A}\) 使得 \(at \in \mathbf{S}\)（换言之，\(\mathbf{T}\) 的每个元素整除 \(\mathbf{S}\) 的某个元素）。
\item
  每个与 T 相交的素理想也与 S 相交。
\end{enumerate}

上面已经看到 \(i_{\mathbf{A}}^{\mathrm{T},\mathrm{S}} = 1_{\mathbb{M}} \otimes i_{\mathbf{A}}^{\mathrm{T},\mathrm{S}}\)，由此立即可证 (a) 与 (b) 等价。令 \(\mathrm{T}' = i_{\mathbf{A}}^{\mathrm{S}}(\mathrm{T})\)；于是（命题 7）\(\mathrm{T}^{-1}\mathbf{A}\) 与 \(\mathrm{T}'^{-1}(\mathrm{S}^{-1}\mathbf{A})\) 等同，而 (a) 等价于说 \(\mathrm{T}'\) 的元素在 \(\mathrm{S}^{-1}\mathbf{A}\) 中可逆（第 1 节，注 5）。现在，说 \((t/1)(a/s) = 1/1\)（\(t \in \mathrm{T}\)、\(a \in \mathrm{A}\)、\(s \in \mathrm{S}\)）意味着存在 \(s' \in \mathrm{S}\) 使得 \(tas' = ss'\)，这就表明 (a) 与 (c) 等价。我们证明 (d) 蕴含 (c)。设 \(t\) 是 \(\mathrm{T}\) 的一个元素，并设 \(t/1\) 在 \(\mathrm{S}^{-1}\mathbf{A}\) 中不可逆；于是存在一个极大理想 \(m'\)（\(\mathrm{S}^{-1}\mathbf{A}\) 的），它包含 \(t/1\)（《代数学》第 I 章，§ 8，第 7 节，定理 2），而 \(p = (i_{\mathbf{A}}^{\mathrm{S}})^{-1}(m')\) 是 \(\mathbf{A}\) 的一个素理想，它包含 \(t\) 且不与 \(\mathbf{S}\) 相交（因为在 \(i_{\mathbf{A}}^{\mathrm{S}}\) 之下 \(\mathbf{S}\) 的元素的像是可逆的）。反之，若存在素理想 \(p\) 与 \(\mathrm{T}\) 相交而不与 \(\mathrm{S}\) 相交，则 \(p \cap \mathrm{T}\) 中没有元素能整除 \(\mathrm{S}\) 的元素；这就证明 (c) 蕴含 (d)，并完成证明。

由命题 8 可知，在 A 的包含 S 且满足命题 8 的诸等价条件的乘性子集 T 当中，存在一个最大的，它由 A 中整除 S 的某个元素的全体元素组成（参见习题 1）。

命题 9. 设 I 是一个右定向的预序集，\((\mathrm{S}_{\alpha})_{\alpha \in \mathbf{I}}\) 是环 A 的乘性子集的一个递增族，而 \(S = \bigcup_{\alpha \in I} S_{\alpha}\)。我们记 \(\rho_{\beta \alpha} = i_{A}^{S_{\beta}, S_{\alpha}}\)（对 \(\alpha \leqslant \beta\)），并记 \(\rho_{\alpha} = i_{A}^{S, S_{\alpha}}\)。则 \((S_{\alpha}^{-1}A, \rho_{\beta \alpha})\) 是一个环的正向系，并且，若对所有 \(\alpha \in I\)，\(\rho_{\alpha}'\) 是 \(\mathbf{S}_{\alpha}^{-1}\mathbf{A}\) 到 \(\lim_{\longrightarrow}\mathbf{S}_{\alpha}^{-1}\mathbf{A}\) 的典范映射，则存在唯一的同构 \(j\)（从 \(\lim_{\longrightarrow}\mathbf{S}_{\alpha}^{-1}\mathbf{A}\) 到 \(\mathbf{S}^{-1}\mathbf{A}\)），使得 \(j\circ \rho_{\alpha}^{\prime} = \rho_{\alpha}\) 对所有 \(\alpha \in \mathbf{I}\) 成立。

对 \(\alpha \leqslant \beta \leqslant \gamma\)，\(\rho_{\gamma \alpha} = \rho_{\gamma \beta} \circ \rho_{\beta \alpha}\)（第 1 节，命题 2 的推论 3），因此 \((S_{\alpha}^{-1}A, \rho_{\beta \alpha})\) 是一个正向系。我们记 \(A' = \varinjlim S_{\alpha}^{-1}A\)；由于 \(\rho_{\alpha} = \rho_{\beta} \circ \rho_{\beta \alpha}\) 对 \(\alpha \leqslant \beta\) 成立（第 1 节，命题 2 的推论 3），\((\rho_{\alpha})\) 是同态的一个正向系，而 \(j = \varinjlim \rho_{\alpha}\) 是从 \(A'\) 到 \(S^{-1}A\) 的唯一同态，使得 \(j \circ \rho_{\alpha}' = \rho_{\alpha}\) 对所有 \(\alpha \in I\) 成立。同态 \(\rho_{\alpha}' \circ i_{A^{\alpha}}^{S}: A \to A'\) 都相等，因为 \(\rho_{\beta \alpha} \circ i_{A^{\alpha}}^{S} = i_{A^{\beta}}^{S}\) 对 \(\alpha \leqslant \beta\) 成立；设 \(u\) 是它们的公共值。显然 \(u(S)\) 的元素在 \(A'\) 中可逆，这表明存在同态 \(h: S^{-1}A \to A'\) 使得 \(h \circ i_{A}^{S} = u\)（第 1 节，命题 1）。则

\[
j \circ h \circ i _ {\mathrm{S}} ^ {\mathrm{A}} = j \circ u = j \circ \rho_ {\alpha} ^ {\prime} \circ i _ {\mathrm{S} ^ {\alpha}} ^ {\mathrm{A}} = \rho_ {\alpha} \circ i _ {\mathrm{A} ^ {\alpha}} ^ {\mathrm{S}} = i _ {\mathrm{A}} ^ {\mathrm{S}}
\]

对所有 \(\alpha \in \mathbf{I}\) 成立，因而 \(j \circ h\) 是 \(\mathbf{S}^{-1}\mathbf{A}\) 的恒等自同构。另一方面，对所有 \(\alpha \in \mathbf{I}\)，

\[
h \circ j \circ \rho_ {\alpha} ^ {\prime} \circ i _ {A ^ {\alpha}} ^ {S} = h \circ \rho_ {\alpha} \circ i _ {A ^ {\alpha}} ^ {S} = h \circ i _ {A} ^ {S} = u = \rho_ {\alpha} ^ {\prime} \circ i _ {A ^ {\alpha}} ^ {S},
\]

从而 \(h \circ j \circ \rho_{\alpha}' = \rho_{\alpha}'\) 对所有 \(\alpha \in I\) 成立；由此得出 \(h \circ j\) 是 \(A'\) 的恒等自同构，因而 \(j\) 是同构。

推论. 在命题 9 的假设下，设 M 是 A-模。我们记 \(f_{\beta \alpha} = i_{\mathrm{M}}^{\mathrm{S}_{\beta},\mathrm{S}_{\alpha}}\) 对 \(\alpha \leqslant \beta, f_{\alpha} = i_{\mathrm{M}}^{\mathrm{S},\mathrm{S}_{\alpha}}\)（对所有 \(\alpha \in \mathbf{I}\)），并设 \(f_{\alpha}'\) 是 \(\mathrm{S}_{\alpha}^{-1}\mathrm{M}\) 到 \(\lim \mathrm{S}_{\alpha}^{-1}\mathrm{M}\) 的典范映射；则存在一个 \(\mathrm{S}^{-1}\mathrm{A}\)-同构 \(g\)（从 \(\mathrm{S}^{-1}\mathrm{M}\) 到 \(\lim \mathrm{S}_{\alpha}^{-1}\mathrm{M}\)），使得 \(g \circ f_{\alpha} = f_{\alpha}'\) 对所有 \(\alpha \in \mathbf{I}\) 成立。

该推论由定义 \(S_{\alpha}^{-1}M = M \otimes_{A} S_{\alpha}^{-1}A\) 与 \(S^{-1}M = M \otimes_{A} S^{-1}A\) 以及取正向极限与张量积交换这一事实（《代数学》第 II 章，§ 6，第 3 节，命题 7）立即得出。

